% Part: first-order-logic
% Chapter: completeness
% Section: maximally-consistent-sets

% Definition of maximally consistent sets. Properties of maximally
% consistent sets required for the completeness proof are in
% maximally-consistent-sets.tex in the chapter on the proof system used.

\documentclass[../../../include/open-logic-section]{subfiles}

\begin{document}

\olfileid{fol}{com}{mcs}
\olsection{مجموعات \usetoken{P}{sentence} المتسقة القصوى}

\begin{defn}[المجموعة المتسقة القصوى]
\ollabel{mcs}
تكون مجموعة~$\Gamma$ من~!!{sentence}s 
\emph{متسقة قصوى} إذا وفقط إذا
\begin{enumerate}
\item كانت~$\Gamma$ متسقة، و
\item إذا كان~$\Gamma \subsetneq \Gamma'$، كانت~$\Gamma'$ غير متسقة.
\end{enumerate}
\end{defn}

\begin{explain}
وفيما يأتي تعريف بديل مكافئ للتعريف السابق: تكون مجموعة~$\Gamma$ من
الجمل \emph{متسقة قصوى} إذا وفقط إذا
\begin{enumerate}
\item كانت~$\Gamma$ متسقة، و
\item إذا كانت~$\Gamma \cup \{ !A \}$ متسقة، كان~$!A \in \Gamma$.
\end{enumerate}
وبعبارة أخرى، لا يمكن إضافة~!!{sentence}s غير الموجودة أصلًا في
$\Gamma$ إلى مجموعة متسقة قصوى~$\Gamma$ من دون جعل المجموعة الأكبر
الناتجة غير متسقة.
\end{explain}

\begin{explain}
للمجموعات المتسقة القصوى أهمية في برهان الاكتمال، لأننا نستطيع ضمان
أن كل مجموعة متسقة من~!!{sentence}s~$\Gamma$ محتواة في مجموعة متسقة
قصوى~$\Gamma^*$، وأن المجموعة المتسقة القصوى تحتوي، لكل~!!{sentence}
$!A$، إما~$!A$ وإما نفيها~$\lnot !A$. ويصدق ذلك بوجه خاص على
!!{sentence}s الذرية؛ ومن ثم نستطيع، انطلاقًا من مجموعة متسقة قصوى في
لغة وُسعت توسيعًا مناسبًا بـ!!{constant}s، أن نبني~!!a{structure}
يُعرّف فيه تأويل~!!{predicate}s بحسب أي~!!{sentence}s ذرية تنتمي إلى
$\Gamma^*$. ويمكن عندئذ إثبات أن هذا~!!{structure} يجعل كل
!!{sentence}s في~$\Gamma^*$، وبالتالي في~$\Gamma$ أيضًا، صادقة. ويتطلب
برهان هذه الحقيقة الأخيرة أن يكون~$\lnot !A \in \Gamma^*$ إذا وفقط
إذا كان~$!A \notin \Gamma^*$، وأن يكون
$(!A \lor !B) \in \Gamma^*$ إذا وفقط إذا كان~$!A \in \Gamma^*$ أو
$!B \in \Gamma^*$، وهلم جرًا.
\end{explain}

\begin{prop}
\ollabel{prop:mcs}
افترض أن~$\Gamma$ متسقة قصوى. عندئذ:
\begin{enumerate}
\item \ollabel{prop:mcs-prov-in} إذا كان~$\Gamma \Proves !A$، كان
$!A \in \Gamma$.

\item \ollabel{prop:mcs-either-or} لكل~$!A$، إما~$!A \in \Gamma$ أو
$\lnot !A \in \Gamma$.

\tagitem{prvAnd}{\ollabel{prop:mcs-and} يكون
$(!A \land !B) \in \Gamma$ إذا وفقط إذا كان كل من
$!A \in \Gamma$ و$!B \in \Gamma$.}{}

\tagitem{prvOr}{\ollabel{prop:mcs-or} يكون
$(!A \lor !B) \in \Gamma$ إذا وفقط إذا كان إما~$!A \in \Gamma$
وإما~$!B \in \Gamma$.}{}

\tagitem{prvIf}{\ollabel{prop:mcs-if} يكون
$(!A \lif !B) \in \Gamma$ إذا وفقط إذا كان إما~$!A \notin \Gamma$
وإما~$!B \in \Gamma$.}{}
\end{enumerate}
\end{prop}

\begin{proof}
لنفترض في كل ما يأتي أن~$\Gamma$ متسقة قصوى.
\begin{enumerate}
\item إذا كان~$\Gamma \Proves !A$، كان~$!A \in \Gamma$.

افترض~$\Gamma \Proves !A$. وافترض، طلبًا للتناقض، أن
$!A \notin \Gamma$. لما كانت~$\Gamma$ متسقة قصوى، تكون
$\Gamma \cup \{!A\}$ غير متسقة، ومن ثم
$\Gamma \cup \{!A\} \Proves \lfalse$. وبموجب
\tagrefs{prfSC/{fol:seq:prv:prop:provability-contr},prfND/{fol:ntd:prv:prop:provability-contr}}،
تكون~$\Gamma$ غير متسقة. وهذا يناقض افتراض اتساقها. لذلك لا يمكن أن
يكون~$!A \notin \Gamma$، فيكون~$!A \in \Gamma$.

\item لكل~$!A$، إما~$!A \in \Gamma$ أو~$\lnot !A \in \Gamma$.

افترض، طلبًا للتناقض، أنه يوجد~$!A$ بحيث~$!A \notin \Gamma$ و
$\lnot !A \notin \Gamma$. ولأن~$\Gamma$ متسقة قصوى، تكون
$\Gamma \cup \{!A\}$ و$\Gamma \cup \{\lnot !A\}$ كلتاهما غير
متسقة، فيكون~$\Gamma \cup \{!A\} \Proves \lfalse$ و
$\Gamma \cup \{\lnot !A\} \Proves \lfalse$. وبموجب
\tagrefs{prfSC/{fol:seq:prv:prop:provability-exhaustive},prfND/{fol:ntd:prv:prop:provability-exhaustive}}،
تكون~$\Gamma$ غير متسقة، وهذا تناقض. ومن ثم لا يوجد~!!a{sentence}
$!A$ كهذا، ولكل~$!A$ إما~$!A \in \Gamma$ أو
$\lnot !A \in \Gamma$.

\tagitem{defAnd}{}{%
\iftag{probAnd}{تمرين.}{%
يكون~$(!A \land !B) \in \Gamma$ إذا وفقط إذا كان كل من
$!A \in \Gamma$ و$!B \in \Gamma$:

في الاتجاه الأمامي، افترض~$(!A \land !B) \in \Gamma$. عندئذ
$\Gamma \Proves !A \land !B$. وبموجب
\tagrefs{prfSC/{fol:seq:prv:prop:provability-land-left},prfND/{fol:ntd:prv:prop:provability-land-left}}،
يكون~$\Gamma \Proves !A$ و$\Gamma \Proves !B$. وبموجب
\olref{prop:mcs-prov-in}، يكون~$!A \in \Gamma$ و$!B \in \Gamma$ كما
أردنا.

وفي الاتجاه العكسي، ليكن~$!A \in \Gamma$ و$!B \in \Gamma$. عندئذ
$\Gamma \Proves !A$ و$\Gamma \Proves !B$. وبموجب
\tagrefs{prfSC/{fol:seq:prv:prop:provability-land-right},prfND/{fol:ntd:prv:prop:provability-land-right}}،
يكون~$\Gamma \Proves !A \land !B$. وبموجب~\olref{prop:mcs-prov-in}،
يكون~$(!A \land !B) \in \Gamma$.}}

\tagitem{defOr}{}{%
\iftag{probOr}{تمرين.}{%
يكون~$(!A \lor !B) \in \Gamma$ إذا وفقط إذا كان إما
$!A \in \Gamma$ وإما~$!B \in \Gamma$.

لإثبات مقابل الاتجاه الأمامي، افترض~$!A \notin \Gamma$ و
$!B \notin \Gamma$. نريد أن نبيّن~$(!A \lor !B) \notin \Gamma$.
ولأن~$\Gamma$ متسقة قصوى، يكون
$\Gamma \cup \{!A\} \Proves \lfalse$ و
$\Gamma \cup \{!B\} \Proves \lfalse$. وبموجب
\tagrefs{prfSC/{fol:seq:prv:prop:provability-lor-left},prfND/{fol:ntd:prv:prop:provability-lor-left}}،
تكون~$\Gamma \cup \{(!A \lor !B)\}$ غير متسقة. ومن ثم
$(!A \lor !B) \notin \Gamma$ كما أردنا.

وفي الاتجاه العكسي، افترض~$!A \in \Gamma$ أو~$!B \in \Gamma$.
عندئذ~$\Gamma \Proves !A$ أو~$\Gamma \Proves !B$. وبموجب
\tagrefs{prfSC/{fol:seq:prv:prop:provability-lor-right},prfND/{fol:ntd:prv:prop:provability-lor-right}}،
يكون~$\Gamma \Proves !A \lor !B$. وبموجب~\olref{prop:mcs-prov-in}،
يكون~$(!A \lor !B) \in \Gamma$ كما أردنا.}}

\tagitem{defIf}{}{%
\iftag{probIf}{تمرين.}{%
يكون~$(!A \lif !B) \in \Gamma$ إذا وفقط إذا كان إما
$!A \notin \Gamma$ وإما~$!B \in \Gamma$:

في الاتجاه الأمامي، ليكن~$(!A \lif !B) \in \Gamma$، وافترض طلبًا
للتناقض أن~$!A \in \Gamma$ و$!B \notin \Gamma$. وبموجب هذين
الافتراضين، يكون~$\Gamma \Proves !A \lif !B$ و
$\Gamma \Proves !A$. وبموجب
\tagrefs{prfSC/{fol:seq:prv:prop:provability-mp},prfND/{fol:ntd:prv:prop:provability-mp}}،
يكون~$\Gamma \Proves !B$. لكن~\olref{prop:mcs-prov-in} تعطي عندئذ
$!B \in \Gamma$، وهو ما يناقض~$!B \notin \Gamma$.

وفي الاتجاه العكسي، تأمل أولًا حالة~$!A \notin \Gamma$. بموجب
\olref{prop:mcs-either-or}، يكون~$\lnot !A \in \Gamma$، ومن ثم
$\Gamma \Proves \lnot !A$. وبموجب
\tagrefs{prfSC/{fol:seq:prv:prop:provability-lif},prfND/{fol:ntd:prv:prop:provability-lif}}،
يكون~$\Gamma \Proves !A \lif !B$. وتطبيق~\olref{prop:mcs-prov-in}
مرة أخرى يعطينا~$(!A \lif !B) \in \Gamma$ كما أردنا.

والآن تأمل حالة~$!B \in \Gamma$. عندئذ~$\Gamma \Proves !B$، وبموجب
\tagrefs{prfSC/{fol:seq:prv:prop:provability-lif},prfND/{fol:ntd:prv:prop:provability-lif}}،
يكون~$\Gamma \Proves !A \lif !B$. وبموجب~\olref{prop:mcs-prov-in}،
يكون~$(!A \lif !B) \in \Gamma$.}}
\end{enumerate}
\end{proof}

\begin{prob}
أكمل برهان~\olref[fol][com][mcs]{prop:mcs}.
\end{prob}

\end{document}
